% ============================================================
%  INTRODUCTION GÉNÉRALE
% ============================================================
\chapter*{Introduction générale}
\addcontentsline{toc}{chapter}{Introduction générale}
\markboth{Introduction générale}{Introduction générale}

\hspace{1cm}Dans un monde de plus en plus connecté, les menaces cybernétiques ne cessent de croître en complexité et en fréquence. Face à cette réalité, l'intelligence artificielle s'impose comme un levier incontournable pour renforcer les dispositifs de défense des systèmes d'information.

Le présent travail s'inscrit dans le cadre du Projet de Fin d'Études pour l'obtention du diplôme d'Ingénieur d'État en \textbf{Intelligence Artificielle et Cybersécurité (IACS)} à l'École Nationale des Sciences Appliquées de Béni Mellal.

\vspace{0.5cm}
La problématique abordée concerne \textit{[décrire brièvement la problématique]}. Pour y répondre, nous proposons une approche combinant les techniques d'apprentissage profond avec les mécanismes de sécurité réseau.

\vspace{0.5cm}
Ce rapport est organisé en trois chapitres principaux :

\begin{itemize}
  \item \textbf{Chapitre I -- Contexte général et problématique :} présentation de l'organisme d'accueil, de la problématique et du cahier des charges.
  \item \textbf{Chapitre II -- État de l'art :} revue de la littérature sur les techniques d'\gls{ia} appliquées à la cybersécurité et analyse comparative des solutions existantes.
  \item \textbf{Chapitre III -- Conception et réalisation :} description de l'architecture proposée, des choix technologiques, de l'implémentation et des résultats obtenus.
\end{itemize}
